\documentclass{article}
\usepackage{amsmath}

\title{On the Convergance of Iteratve Methods}
\author{Test Author}
\date{}

\begin{document}
\maketitle

\section{Introduction}
This paper studys the convergance of iterative methods for solving linar systems.
The importent contribution of this work is that we provides a new bound on the convergence rate.
Previous work have shown that such methods converges under certain conditions,
but our result are more tighter and applies to a broader class of problems.

\section{Preliminaries}
Let $A$ be a $n \times n$ matrix. We assume that $A$ is positive definite,
which means that all its eigenvalues is positive.
The spectral radius of a matrix $B$, denoted $\rho(B)$, is the largest absolute value of it's eigenvalues.

We recall that a iterative method of the form
\[
x^{(k+1)} = Bx^{(k)} + c
\]
converges if and only if $\rho(B) < 1$. This is a well-known results.

\section{Main Result}
\begin{theorem}
Let $A$ be a symmetric positive definite matrix with eigenvalues $0 < \lambda_1 \le \cdots \le \lambda_n$.
Then the Jacobi method converges and the convergence rate satisfyes
\[
\rho(B_J) \le 1 - \frac{2\lambda_1}{\lambda_1 + \lambda_n}.
\]
\end{theorem}

\begin{proof}
The proof follows from standard spectral analysis. We ommit the details.
\end{proof}

\section{Conclusion}
We have showed that the Jacobi method converges for symmetric positive definite matrices.
The bound we derived are sharp and cannot be improved further without additional assumptions.
This result has importent implications for practical computations.

\end{document}
